% openGauss安装文档
% openGauss安装手册.tex

\documentclass[12pt,UTF8]{ctexbook}

% 设置纸张信息。
\input{../../../../Book/common/page}

% 设置字体，并解决显示难检字问题。
\xeCJKsetup{AutoFallBack=true}
% 注意：ctexbook类已默认设置SimSun为CJKrmdefault
% 以下设置用于确保字体回退和扩展字体可用
% 仅设置扩展字体以避免与默认设置冲突
\setCJKfamilyfont{hei}{SimHei}
\setCJKfamilyfont{kai}{KaiTi}

% 目录 chapter 级别加点（.）。
\usepackage{titletoc}
\titlecontents{chapter}[0pt]{\vspace{3mm}\bf\addvspace{2pt}\filright}{\contentspush{\thecontentslabel\hspace{0.8em}}}{}{\titlerule*[8pt]{.}\contentspage}

% 支持目录点击跳转
\usepackage[colorlinks,linkcolor=blue,citecolor=blue,urlcolor=blue]{hyperref}

% 设置 part 和 chapter 标题格式。
\ctexset{
	part/name= {第,卷},
	part/number={\chinese{part}},
	chapter/name={第,篇},
	chapter/number={\arabic{chapter}}
}

% 图片相关设置。
\usepackage{graphicx}
\graphicspath{{Images/}}

% 设置署名格式。
\newenvironment{shuming}{\hfill\zihao{4}}

% 注脚每页重新编号，避免编号过大。
\usepackage[perpage]{footmisc}

% 设置段落间距
\setlength{\parskip}{0.8em plus 0.2em minus 0.1em}

% 列表格式支持，列表项向右偏移。
\usepackage{enumitem}

% 代码块支持
\usepackage{listings}
\usepackage{xcolor}
\lstset{
    basicstyle=\ttfamily\footnotesize,
    keywordstyle=\color{blue},
    commentstyle=\color{green!60!black},
    stringstyle=\color{red},
    breaklines=true,
    frame=single,
    numbers=left,
    numberstyle=\tiny\color{gray},
    captionpos=b,
    tabsize=4,
    literate={_}{{\char`\_}}1
}

\title{\heiti\zihao{0} openGauss安装手册}
\author{WangFei}
\date{\today}

\begin{document}

\maketitle
\tableofcontents

\frontmatter
\chapter{前言}

openGauss是一款开源关系型数据库管理系统，采用木兰宽松许可证v2发行。openGauss内核源自PostgreSQL，深度融合华为在数据库领域多年的经验，结合企业级场景需求，在架构、事务、存储引擎、优化器等众多领域进行了大量的改进和优化。

本手册旨在帮助用户快速掌握openGauss的安装、配置和使用方法。

\mainmatter

\chapter{openGauss简介}

\section{产品概述}

openGauss是华为开源的关系型数据库管理系统，具有以下特点：

\begin{itemize}[leftmargin=2cm]
    \item \textbf{高性能}：支持多核架构、高性能存储引擎
    \item \textbf{高可用}：支持主备复制、故障自动切换
    \item \textbf{高安全}：支持全密态、透明数据加密
    \item \textbf{易运维}：提供丰富的监控和管理工具
    \item \textbf{开源开放}：完全开源，社区活跃
\end{itemize}

\section{系统架构}

openGauss采用客户端/服务器架构，主要包含以下组件：

\begin{itemize}[leftmargin=2cm]
    \item \textbf{客户端}：提供用户接口，发送SQL请求
    \item \textbf{服务器}：处理SQL请求，管理数据
    \item \textbf{存储引擎}：负责数据的存储和检索
    \item \textbf{优化器}：生成最优的执行计划
\end{itemize}

\chapter{安装前准备}

\section{硬件要求}

准备软硬件安装环境

本章节描述安装前需要进行的环境准备。建议部署openGauss的各服务器具有等价的软硬件配置。

\subsection{硬件环境要求}

表1 硬件环境要求列出了openGauss服务器应具备的最低硬件要求。在实际产品中，硬件配置的规划需考虑数据规模及所期望的数据库响应速度。请根据实际情况进行规划。

\begin{table}[htbp]
    \centering
    \caption{硬件环境要求}
    \label{tab:hardware_requirements}
    \begin{tabular}{|l|p{10cm}|}
    \hline
    \textbf{项目} & \textbf{配置描述} \\ \hline
    内存 & 功能调试建议32GB以上。\\ \hline
    & 性能测试和商业部署时，单实例部署建议128GB以上。\\ \hline
    & 复杂的查询对内存的需求量比较高，在高并发场景下，可能出现内存不足。此时建议使用大内存的机器，或使用负载管理限制系统的并发。\\ \hline
    CPU & 功能调试最小1×8 核 2.0GHz。\\ \hline
    & 性能测试和商业部署时，建议1×16核 2.0GHz。\\ \hline
    & CPU超线程和非超线程两种模式都支持。\\ \hline
    & 说明：个人开发者最低配置2核4G, 推荐配置4核8G。\\ \hline
    & 目前，openGauss仅支持ARM服务器和基于\texttt{X86\_64}通用PC服务器的CPU。\\ \hline
    硬盘 & 用于安装openGauss的硬盘需最少满足如下要求：\\ \hline
    & 至少1GB用于安装openGauss的应用程序。\\ \hline
    & 每个主机需大约300MB用于元数据存储。\\ \hline
    & 预留70\%以上的磁盘剩余空间用于数据存储。\\ \hline
    & 建议系统盘配置为Raid1，数据盘配置为Raid5，且规划4组Raid5数据盘用于安装openGauss。有关Raid的配置方法在本手册中不做介绍。\\ \hline
    & 请参考硬件厂家的手册或互联网上的方法进行配置，其中Disk Cache Policy一项需要设置为Disabled，否则机器异常掉电后有数据丢失的风险。\\ \hline
    & openGauss支持使用SSD盘作为数据库的主存储设备，支持SAS接口和NVME协议的SSD盘，以RAID的方式部署使用。\\ \hline
    网络 & 300兆以上以太网。\\ \hline
    & 建议网卡设置为双网卡冗余bond。有关网卡冗余bond的配置方法在本手册中不做介绍。\\ \hline
    & 请参考硬件厂商的手册或互联网上的方法进行配置。\\ \hline
    \end{tabular}
\end{table}

\subsection{客户端配置}

\begin{itemize}[leftmargin=2cm]
    \item CPU：双核及以上
    \item 内存：2GB及以上
    \item 硬盘：10GB及以上
\end{itemize}

\section{软件要求}

\subsection{操作系统支持}

openGauss支持以下操作系统：

\begin{table}[htbp]
    \centering
    \caption{软件环境要求}
    \label{tab:software_requirements}
    \begin{tabular}{|l|l|}
    \hline
    \textbf{软件类型} & \textbf{配置描述} \\ \hline
    Linux操作系统（ARM） & openEuler 20.03 LTS（推荐采用此操作系统）\\ \hline
    & openEuler 22.03 LTS \\ \hline
    & 麒麟V10 \\ \hline
    & Asianux 7.5 \\ \hline
    & 统信V20 \\ \hline
    Linux操作系统（X86） & openEuler 20.03 LTS \\ \hline
    & openEuler 22.03 LTS \\ \hline
    & CentOS 7.6 \\ \hline
    & Asianux 7.6 \\ \hline
    & 麒麟V10 \\ \hline
    Linux文件系统 & 剩余inode个数 > 15亿（推荐）\\ \hline
    工具 & bzip2 \\ \hline
    \end{tabular}
\end{table}

\textbf{说明}：
\begin{itemize}[leftmargin=2cm]
    \item 当前安装包只能在英文操作系统上安装使用；
    \item OM工具已经支持对基于openEuler/Centos等商业操作系统的安装使用，具体配置信息可以查看OM中的osid.conf文件。
\end{itemize}

\subsection{依赖软件包}

openGauss的软件依赖要求如下表所示：

\begin{table}[htbp]
    \centering
    \caption{软件依赖要求}
    \label{tab:dependency_requirements}
    \begin{tabular}{|l|l|}
    \hline
    \textbf{所需软件} & \textbf{建议版本} \\ \hline
    libaio-devel & 建议版本：0.3.109-13 \\ \hline
    readline-devel & 建议版本：7.0-13 \\ \hline
    libnsl（openEuler+x86环境中） & 建议版本：2.28-36 \\ \hline
    \end{tabular}
\end{table}

\textbf{说明}：建议使用上述操作系统安装光盘或者源中的默认安装包，若不存在上述软件，可参看软件对应的建议版本。

使用yum命令安装依赖包：

\begin{lstlisting}[language=bash, caption=安装依赖软件包]
yum install libaio-devel readline-devel libnsl -y
\end{lstlisting}

\textbf{说明}：
\begin{itemize}[leftmargin=2cm]
    \item 使用\texttt{-y}参数可以自动确认安装，无需手动交互
    \item 如果是其他依赖包（如flex、bison等），也可以通过类似方式安装
\end{itemize}

\section{网络配置}

\subsection{端口规划}

openGauss默认使用以下端口：

\begin{table}[htbp]
    \centering
    \caption{openGauss端口配置}
    \label{tab:port_config}
    \begin{tabular}{|l|l|l|}
    \hline
    \textbf{端口} & \textbf{用途} & \textbf{说明} \\ \hline
    8000 & 数据库服务端口 & 客户端连接端口 \\ \hline
    8001 & HA端口 & 主备复制端口 \\ \hline
    8002 & CM端口 & 集群管理端口 \\ \hline
    \end{tabular}
\end{table}

\subsection{SELinux配置}

修改/etc/selinux/config文件中的"SELINUX"值为"disabled"：

\begin{lstlisting}[language=bash, caption=修改SELinux配置]
vim /etc/selinux/config
\end{lstlisting}

打开文件后，修改"SELINUX"的值为"disabled"：

\begin{lstlisting}[language=bash, caption=SELinux配置内容]
SELINUX=disabled
\end{lstlisting}

执行:wq保存并退出修改，然后重新启动操作系统：

\begin{lstlisting}[language=bash, caption=重启操作系统]
reboot
\end{lstlisting}

\subsection{防火墙配置}

需要开放相关端口，例如：

\begin{lstlisting}[language=bash, caption=开放防火墙端口]
firewall-cmd --zone=public --add-port=8000/tcp --permanent
firewall-cmd --zone=public --add-port=8001/tcp --permanent
firewall-cmd --reload
\end{lstlisting}

或者直接关闭防火墙（适用于测试环境）：

\begin{lstlisting}[language=bash, caption=关闭防火墙]
systemctl stop firewalld
systemctl disable firewalld
\end{lstlisting}

\subsection{关闭RemoveIPC}

在各数据库节点上，关闭RemoveIPC。CentOS操作系统无该参数，可以跳过该步骤。

\textbf{操作步骤：}

修改/etc/systemd/logind.conf文件中的"RemoveIPC"值为"no"：

\begin{lstlisting}[language=bash, caption=修改logind.conf]
vim /etc/systemd/logind.conf
\end{lstlisting}

修改"RemoveIPC"值为"no"：

\begin{lstlisting}[language=bash, caption=RemoveIPC配置]
RemoveIPC=no
\end{lstlisting}

修改/usr/lib/systemd/system/systemd-logind.service文件中的"RemoveIPC"值为"no"：

\begin{lstlisting}[language=bash, caption=修改systemd-logind.service]
vim /usr/lib/systemd/system/systemd-logind.service
\end{lstlisting}

修改"RemoveIPC"值为"no"：

\begin{lstlisting}[language=bash, caption=RemoveIPC配置]
RemoveIPC=no
\end{lstlisting}

重新加载配置参数：

\begin{lstlisting}[language=bash, caption=重新加载配置]
systemctl daemon-reload
systemctl restart systemd-logind
\end{lstlisting}

检查修改是否生效：

\begin{lstlisting}[language=bash, caption=检查配置]
loginctl show-session | grep RemoveIPC
systemctl show systemd-logind | grep RemoveIPC
\end{lstlisting}

\section{用户和目录准备}

\subsection{创建用户组}

创建用户组dbgroup：

\begin{lstlisting}[language=bash, caption=创建用户组]
groupadd dbgroup
\end{lstlisting}

\subsection{创建用户}

创建用户组dbgroup下的普通用户omm，并设置普通用户omm的密码，密码建议设置为44PMLO0a：

\begin{lstlisting}[language=bash, caption=创建数据库用户]
useradd -g dbgroup -m -d /home/omm omm
passwd omm
\end{lstlisting}

\textbf{说明}：
\begin{itemize}[leftmargin=2cm]
    \item \texttt{-g dbgroup}：指定用户所属的用户组
    \item \texttt{-m}：自动创建用户的主目录
    \item \texttt{-d /home/omm}：指定用户的主目录路径
    \item 设置密码时建议使用复杂密码（至少8位），如\texttt{44PMLO0a}
\end{itemize}

\chapter{openGauss安装}

\section{获取安装包}

\subsection{官方下载}

访问openGauss官方网站下载安装包：

\begin{itemize}[leftmargin=2cm]
    \item 网址：https://opengauss.org/zh/download/
    \item 选择对应的版本和操作系统
    \item 下载企业版或轻量版安装包
\end{itemize}

\subsection{推荐下载链接}

openGauss 6.0.5 版本（openEuler 20.03 LTS x86\_64）：

\begin{itemize}[leftmargin=2cm]
    \item \textbf{下载地址}：\url{https://opengauss.obs.cn-south-1.myhuaweicloud.com/6.0.5/openEuler20.03/x86/openGauss-All-6.0.5-openEuler20.03-x86_64.tar.gz}
\end{itemize}

\textbf{说明}：此版本适用于以下操作系统：
\begin{itemize}[leftmargin=2cm]
    \item openEuler 20.03 LTS（x86\_64）
    \item UnionTech UOS 20（x86\_64）
    \item Kylin V10（x86\_64）
\end{itemize}

\subsection{安装包选择说明}

不同操作系统需要选择对应的安装包：

\begin{table}[htbp]
    \centering
    \caption{操作系统与安装包对应关系}
    \label{tab:os_package}
    \begin{tabular}{|l|l|l|}
    \hline
    \textbf{操作系统} & \textbf{架构} & \textbf{对应安装包} \\ \hline
    CentOS 7.6 & x86\_64 & openGauss-6.0.5-CentOS-64bit.tar.gz \\ \hline
    openEuler 20.03 LTS & x86\_64 & openGauss-6.0.5-openEuler-64bit.tar.gz \\ \hline
    openEuler 20.03 LTS & aarch64 & openGauss-6.0.5-openEuler-aarch64.tar.gz \\ \hline
    Kylin V10 & x86\_64 & openGauss-6.0.5-openEuler-64bit.tar.gz \\ \hline
    Kylin V10 & aarch64 & openGauss-6.0.5-openEuler-aarch64.tar.gz \\ \hline
    UnionTech UOS 20 & x86\_64 & openGauss-6.0.5-openEuler-64bit.tar.gz \\ \hline
    UnionTech UOS 20 & aarch64 & openGauss-6.0.5-openEuler-aarch64.tar.gz \\ \hline
    \end{tabular}
\end{table}

\textbf{重要提示}：

\begin{itemize}[leftmargin=2cm]
    \item \textbf{UOS v20（x86\_64）}：需要下载\textbf{openEuler 20.03 LTS x86\_64}版本的安装包
    \item \textbf{Kylin V10（x86\_64）}：同样需要下载openEuler版本的安装包
    \item 推荐下载6.0.5版本，功能更完善
\end{itemize}

\subsection{版本选择}

\begin{table}[htbp]
    \centering
    \caption{openGauss版本对比}
    \label{tab:version_compare}
    \begin{tabular}{|l|l|l|}
    \hline
    \textbf{版本} & \textbf{特点} & \textbf{适用场景} \\ \hline
    企业版 & 功能完整，支持集群 & 生产环境 \\ \hline
    轻量版 & 体积小，安装简单 & 开发测试 \\ \hline
    \end{tabular}
\end{table}

\section{单机安装}

\subsection{解压安装包}

首先在root目录下解压安装包，然后移动到/opt目录：

\begin{lstlisting}[language=bash, caption=解压安装包]
cd /root/
mkdir openGauss
tar -zvxf openGauss-All-6.0.5-openEuler20.03-x86_64.tar.gz
tar jxvf openGauss-Server-6.0.5-openEuler20.03-x86_64.tar.bz2 -C openGauss
ls -lb openGauss/
\end{lstlisting}

\textbf{说明}：
\begin{itemize}[leftmargin=2cm]
    \item \texttt{cd /root/}：切换到root目录
    \item \texttt{tar -zvxf ...tar.gz}：解压外层gzip压缩包，\texttt{-z}表示gzip压缩
    \item \texttt{tar jxvf ...tar.bz2 -C openGauss}：解压内层bzip2压缩包到指定目录
    \item \texttt{ls -lb openGauss/}：查看解压后的文件列表
\end{itemize}

解压完成后，将目录移动到/opt并设置权限：

\begin{lstlisting}[language=bash, caption=移动目录并设置权限]
mv /root/openGauss /opt/
chown -R omm:dbgroup /opt/openGauss
\end{lstlisting}

\subsection{配置依赖库}

openGauss依赖libreadline.so.7，需要在/usr/lib64目录下将libreadline.so.8链接为libreadline.so.7：

\begin{lstlisting}[language=bash, caption=配置libreadline库]
cd /usr/lib64
ln -s libreadline.so.8 libreadline.so.7
\end{lstlisting}

\textbf{说明}：此步骤解决openGauss对旧版本libreadline库的依赖问题。

\subsection{使用simpleInstall安装}

切换到omm用户，然后进入simpleInstall目录：

\begin{lstlisting}[language=bash, caption=切换用户并进入simpleInstall目录]
su - omm
cd /opt/openGauss/simpleInstall
\end{lstlisting}

执行install.sh脚本安装openGauss：

\begin{lstlisting}[language=bash, caption=执行安装脚本]
sh install.sh -w 44PMLO0a
\end{lstlisting}

\textbf{说明}：
\begin{itemize}[leftmargin=2cm]
    \item \texttt{-w}参数用于设置初始化数据库密码（gs\_initdb指定），安全需要必须设置
    \item 密码需满足复杂度要求：至少8位，包含大小写字母、数字和特殊字符
\end{itemize}

\subsection{验证安装}

安装执行完成后，使用ps和gs\_ctl查看进程是否正常：

\begin{lstlisting}[language=bash, caption=Check process status]
ps ux | grep gaussdb
gs_ctl query -D /opt/openGauss/data/single_node
\end{lstlisting}

执行ps命令，显示类似如下信息：

\begin{verbatim}
omm      24209 11.9  1.0 1852000 355816 pts/0  Sl   01:54   0:33 \
/opt/openGauss/bin/gaussdb -D /opt/openGauss/data/single_node
omm      20377  0.0  0.0 119880  1216 pts/0    S+   15:37   0:00 \
grep --color=auto gaussdb
\end{verbatim}

执行gs\_ctl命令，显示类似如下信息：

\begin{verbatim}
gs_ctl query ,datadir is /opt/openGauss/data/single_node
HA state:
    local_role                     : Normal
    static_connections             : 0
    db_state                       : Normal
    detail_information             : Normal

Senders info:
    No information

Receiver info:
    No information
\end{verbatim}

\textbf{状态说明}：
\begin{itemize}[leftmargin=2cm]
    \item \textbf{local\_role: Normal}：本地角色正常
    \item \textbf{db\_state: Normal}：数据库状态正常
    \item \textbf{detail\_information: Normal}：详细信息正常
\end{itemize}

\section{Docker安装}

\subsection{拉取镜像}

\begin{lstlisting}[language=bash, caption=拉取openGauss Docker镜像]
docker pull enmotech/opengauss:latest
# 或指定版本，如：enmotech/opengauss:3.0.0
\end{lstlisting}

\subsection{启动容器}

\begin{lstlisting}[language=bash, caption=启动openGauss容器]
docker run --name opengauss \
  -p 5432:5432 \
  -e GS_PASSWORD=Enmo@1234 \
  -d enmotech/opengauss:latest
\end{lstlisting}

\subsection{连接数据库}

\begin{lstlisting}[language=bash, caption=连接Docker中的数据库]
docker exec -it opengauss gsql -d postgres -p 5432
\end{lstlisting}

\chapter{数据库配置}

\section{初始化配置}

\subsection{连接数据库}

使用gsql命令连接到数据库：

\begin{lstlisting}[language=bash, caption=连接数据库]
gsql -d postgres -p 5432
\end{lstlisting}

\textbf{说明}：
- \texttt{-d}：指定要连接的数据库名
- \texttt{-p}：指定数据库端口号

\subsection{修改密码}

\begin{lstlisting}[language=sql, caption=修改管理员密码]
ALTER USER omm WITH PASSWORD 'NewPassword@123';
\end{lstlisting}

\subsection{创建数据库}

\begin{lstlisting}[language=sql, caption=创建新数据库]
CREATE DATABASE mydb ENCODING 'UTF8';
\end{lstlisting}

\subsection{创建用户}

\begin{lstlisting}[language=sql, caption=创建普通用户]
CREATE USER user1 WITH PASSWORD 'UserPassword@123';
GRANT ALL PRIVILEGES ON DATABASE mydb TO user1;
\end{lstlisting}

\section{参数配置}

\subsection{查看当前参数}

\begin{lstlisting}[language=sql, caption=查看数据库参数]
SHOW ALL;
\end{lstlisting}

\subsection{修改参数}

\begin{lstlisting}[language=sql, caption=修改连接数参数]
ALTER SYSTEM SET max_connections = 200;
\end{lstlisting}

修改参数后需要重启数据库：

\begin{lstlisting}[language=bash, caption=重启数据库]
gs_ctl restart -D /opt/openGauss/data/single_node -Z single_node
\end{lstlisting}

\subsection{常用参数说明}

\begin{table}[htbp]
    \centering
    \caption{常用配置参数}
    \label{tab:common_params}
    \begin{tabular}{|l|l|l|}
    \hline
    \textbf{参数名} & \textbf{默认值} & \textbf{说明} \\ \hline
    max\_connections & 100 & 最大连接数 \\ \hline
    shared\_buffers & 128MB & 共享缓冲区大小 \\ \hline
    effective\_cache\_size & 4GB & 有效缓存大小 \\ \hline
    maintenance\_work\_mem & 16MB & 维护操作内存 \\ \hline
    checkpoint\_completion\_target & 0.5 & 检查点完成目标 \\ \hline
    \end{tabular}
\end{table}

\section{网络配置}

\subsection{配置监听地址}

修改\texttt{/opt/openGauss/data/single\_node/postgresql.conf}文件：

\begin{lstlisting}[language=bash, caption=配置监听地址]
listen_addresses = '*'
port = 8000
\end{lstlisting}

\subsection{配置访问控制}

修改\texttt{/opt/openGauss/data/single\_node/pg\_hba.conf}文件：

\begin{lstlisting}[language=bash, caption=配置访问控制规则]
# 允许所有IP连接
host    all             all             0.0.0.0/0               sha256

# 允许特定网段连接
host    all             all             192.168.1.0/24          sha256
\end{lstlisting}

\chapter{数据库管理}

\section{启动和停止}

\subsection{启动数据库}

\begin{lstlisting}[language=bash, caption=启动数据库]
gs_ctl start -D /opt/openGauss/data/single_node -Z single_node
\end{lstlisting}

\subsection{停止数据库}

\begin{lstlisting}[language=bash, caption=停止数据库]
gs_ctl stop -D /opt/openGauss/data/single_node -Z single_node
\end{lstlisting}

\subsection{重启数据库}

\begin{lstlisting}[language=bash, caption=重启数据库]
gs_ctl restart -D /opt/openGauss/data/single_node -Z single_node
\end{lstlisting}

\subsection{查看状态}

\begin{lstlisting}[language=bash, caption=查看数据库状态]
gs_ctl status -D /opt/openGauss/data/single_node
\end{lstlisting}

\section{备份和恢复}

\subsection{逻辑备份}

\begin{verbatim}
gs_dump mydb -f mydb_backup.sql -U omm
\end{verbatim}

\subsection{逻辑恢复}

\begin{verbatim}
gsql -d mydb -f mydb_backup.sql
\end{verbatim}

\subsection{物理备份}

\begin{lstlisting}[language=bash, caption=物理备份]
gs_basebackup -D /backup/basebackup -Fp -Xs -P -U omm
\end{lstlisting}

\section{日志管理}

\subsection{查看日志位置}

\begin{lstlisting}[language=bash, caption=查看日志目录]
ls -lh /var/log/omm/
\end{lstlisting}

\subsection{查看错误日志}

\begin{lstlisting}[language=bash, caption=查看数据库日志]
tail -f /var/log/omm/postgresql-*.log
\end{lstlisting}

\subsection{日志清理}

\begin{lstlisting}[language=sql, caption=清理旧日志]
SELECT pg_rotate_logfile();
\end{lstlisting}

\chapter{性能优化}

\section{SQL优化}

\subsection{查看执行计划}

\begin{lstlisting}[language=sql, caption=查看SQL执行计划]
EXPLAIN ANALYZE SELECT * FROM table_name WHERE condition;
\end{lstlisting}

\subsection{创建索引}

\begin{lstlisting}[language=sql, caption=创建索引]
CREATE INDEX idx_table_column ON table_name(column_name);
\end{lstlisting}

\subsection{分析表统计信息}

\begin{lstlisting}[language=sql, caption=更新表统计信息]
ANALYZE table_name;
\end{lstlisting}

\section{内存优化}

\subsection{调整共享缓冲区}

\begin{lstlisting}[language=sql, caption=调整共享缓冲区大小]
ALTER SYSTEM SET shared_buffers = '1GB';
\end{lstlisting}

\subsection{调整工作内存}

\begin{lstlisting}[language=sql, caption=调整工作内存]
ALTER SYSTEM SET work_mem = '64MB';
\end{lstlisting}

\section{查询优化}

\subsection{启用慢查询日志}

\begin{lstlisting}[language=sql, caption=启用慢查询日志]
ALTER SYSTEM SET log_min_duration_statement = '1000';
\end{lstlisting}

\subsection{查看慢查询}

\begin{lstlisting}[language=sql, caption=查询慢查询日志]
SELECT * FROM pg_stat_statements ORDER BY total_time DESC LIMIT 10;
\end{lstlisting}

\chapter{安全管理}

\section{用户权限管理}

\subsection{查看用户列表}

\begin{lstlisting}[language=sql, caption=查看所有用户]
SELECT usename FROM pg_user;
\end{lstlisting}

\subsection{授予权限}

\begin{lstlisting}[language=sql, caption=授予用户权限]
GRANT SELECT, INSERT, UPDATE ON table_name TO user_name;
\end{lstlisting}

\subsection{撤销权限}

\begin{lstlisting}[language=sql, caption=撤销用户权限]
REVOKE ALL PRIVILEGES ON table_name FROM user_name;
\end{lstlisting}

\section{数据加密}

\subsection{透明数据加密}

openGauss支持透明数据加密（TDE），保护静态数据安全。

\begin{lstlisting}[language=sql, caption=启用透明数据加密]
CREATE EXTENSION gs_encrypted;
\end{lstlisting}

\subsection{SSL连接配置}

配置SSL加密连接：

\begin{lstlisting}[language=bash, caption=配置SSL]
ssl = on
ssl_cert_file = 'server.crt'
ssl_key_file = 'server.key'
\end{lstlisting}

\section{审计日志}

\subsection{启用审计}

\begin{lstlisting}[language=sql, caption=启用审计功能]
ALTER SYSTEM SET audit_enabled = on;
\end{lstlisting}

\subsection{查看审计日志}

\begin{lstlisting}[language=bash, caption=查看审计日志]
cat /var/log/omm/pg_audit/*.log
\end{lstlisting}

\chapter{故障排查}

\section{常见问题}

\subsection{数据库无法启动}

\begin{itemize}[leftmargin=2cm]
    \item 检查端口是否被占用
    \item 检查磁盘空间是否充足
    \item 检查配置文件是否正确
    \item 查看错误日志获取详细信息
\end{itemize}

\subsection{连接失败}

\begin{itemize}[leftmargin=2cm]
    \item 检查网络连接
    \item 检查防火墙设置
    \item 检查pg\_hba.conf配置
    \item 验证用户名和密码
\end{itemize}

\subsection{性能问题}

\begin{itemize}[leftmargin=2cm]
    \item 检查系统资源使用情况
    \item 分析慢查询日志
    \item 检查锁等待情况
    \item 优化SQL语句
\end{itemize}

\section{日志分析}

\subsection{错误日志分析}

\begin{lstlisting}[language=bash, caption=分析错误日志]
grep "ERROR" /var/log/omm/postgresql-*.log | tail -20
\end{lstlisting}

\subsection{性能日志分析}

\begin{lstlisting}[language=sql, caption=查看性能统计]
SELECT * FROM pg_stat_activity;
\end{lstlisting}

\chapter{附录}

\section{常用命令速查}

\begin{table}[htbp]
    \centering
    \caption{常用命令速查表}
    \label{tab:command_reference}
    \begin{tabular}{|l|l|}
    \hline
    \textbf{命令} & \textbf{说明} \\ \hline
    gs\_ctl start & 启动数据库 \\ \hline
    gs\_ctl stop & 停止数据库 \\ \hline
    gs\_ctl restart & 重启数据库 \\ \hline
    gs\_ctl status & 查看状态 \\ \hline
    gs\_dump & 逻辑备份 \\ \hline
    gs\_restore & 逻辑恢复 \\ \hline
    gsql & 命令行客户端 \\ \hline
    \end{tabular}
\end{table}

\section{参考资料}

\begin{itemize}[leftmargin=2cm]
    \item openGauss官方文档：https://opengauss.org/zh/docs/
    \item openGauss社区：https://opengauss.org/zh/community/
    \item openGauss源码：https://gitee.com/opengauss/openGauss-server
\end{itemize}

\section{技术支持}

\begin{itemize}[leftmargin=2cm]
    \item 官方论坛：https://opengauss.org/zh/forum/
    \item 邮件列表：https://opengauss.org/zh/maillist/
    \item 问题反馈：https://gitee.com/opengauss/openGauss-server/issues
\end{itemize}

\backmatter

\end{document}